\documentclass{article}
\usepackage[T1]{fontenc}
\usepackage{lmodern}
\usepackage{graphicx}
\usepackage{amsmath,amssymb}
\usepackage[a4paper,margin=2cm]{geometry}
\usepackage[absolute,overlay]{textpos}
\setlength{\TPHorizModule}{1mm}
\setlength{\TPVertModule}{1mm}
\usepackage[indonesian]{babel}
\usepackage{hyperref}
\setlength{\emergencystretch}{2em}
% unit: Halaman 7

\begin{document}

% === Halaman 7 (python/native) ===

• Hasil penjumlahan Ri–1 dengan Ki menghasilkan 

luaran yang panjangnya 32 bit. 

• Luaran ini dibagi mejadi 8 bagian yang masing-

masing panjangnya 4 bit. 

• Setiap 4 bit masuk ke dalam kotak S untuk proses 

substitusi. Empat bit pertama masuk ke dalam 

kotak S pertama, 4 bit kedua masuk ke dalam kotak 

S kedua, demikian seterusnya. 

• Hasil substitusi setiap kotak S adalah 4 bit. GOST 

memiliki 8 buah kotak S, setiap kotak berisi 16 

buah elemen nilai. Setiap kotak berisi permutasi 

angka 0 sampai 15. 

7

Li - 1

Ri - 1

Substitusi

Kotak-S

Pergeseran sirkuler

ke kiri 11 bit



Ri

Li

Ki

Keterangan:

adalah operator penjumlahan dalam modulo 232

f

\end{document}
